\section{Introduction}
\label{sec:intro}
Retrieval-Augmented Generation (RAG) grounds the answers of a Large Language Model (LLM) in retrieved evidence \cite{NEURIPS2020_6b493230}. In practice the evidence rarely arrives all at once. A retriever returns passages of mixed relevance, an agent fetches them over several rounds, and for multi-hop questions the answer depends on several separate facts that may be retrieved at different times \cite{trivedi2022musique}. Two questions then matter at every step. Is the evidence now \emph{sufficient} to answer? And, if it is, will the model actually \emph{take it up}, or will it ignore the decisive passage and answer as before? Today both questions are usually put to the model itself: it is told to say ``INSUFFICIENT'' when it cannot answer, and its output is taken as the verdict.

This paper asks whether the model's internal state gives a better answer. Prior work has shown that residual-stream activations carry information about the conflict between retrieved context and parametric memory, and about how much the model relies on the context \cite{SpARE,context-vs-memory}. Nothing comparable exists for sufficiency and uptake, and we know of no study that compares an internal readout of these properties with the model's own verbalised judgement, or with a text classifier given the same input.

We study this with \emph{evidence trajectories}. For every multi-hop question we add the gold paragraphs one at a time and run a fresh forward pass on each context, reading the residual stream at the position that produces the first answer token, before anything is generated. The addition that makes the evidence complete, the \emph{decisive increment}, is paired with controls that add, to the same starting context, a distractor of the same length, a duplicate of a paragraph already present, a falsified version of the missing paragraph, or an irrelevant paragraph that mentions the answer string. A probe trained to tell the decisive increment from these controls cannot rely on length, on the presence of the answer string or on lexical relevance alone. Every increment is also labelled with what the model's answer did, so a second probe, trained on decisive increments only, predicts uptake: whether the answer becomes correct, or the new evidence is ignored.

Our findings are, in order of importance, the following. First, sufficiency and uptake are both linearly readable before generation, at $0.95$ and $0.78$ AUROC, on splits that keep supporting documents, answer entities and relation types disjoint between training and test. Second, these readouts compare favourably with the alternatives. For sufficiency the internal state is far more accurate than the model's verbalised judgement, which declares complete evidence insufficient a third of the time, and somewhat more accurate than a text classifier that sees the whole context. For uptake no text classifier gets far above chance; the model's decision to answer at all predicts it about as well as the probe does, and the probe adds to that decision and is the best single predictor of whether an answer the model gives is correct. Third, as a corollary, the probe directions are not levers: steering, patching and projecting along them leave behaviour unchanged. The one direction with a causal effect is the difference between the mean sufficient and insufficient states, and it acts as a gate on abstention rather than on evidence use, which connects our results to work on steering abstention through confidence directions \cite{ActivationSteering,Kumaran2026}. Fourth, the readouts are recoverable in three model families and on other datasets, and directions learned on \musique{} transfer to HotpotQA almost intact but only partially to RAGBench's naturally retrieved documents and to synthetic evidence.

Our contributions are as follows.
\begin{itemize}
  \item \textbf{An increment-based protocol} for reading evidence sufficiency and evidence uptake from a single forward pass before generation, with matched controls that rule out the surface explanations.
  \item \textbf{A comparison with the baselines that matter}: text classifiers on the same inputs, and the model's own verbalised judgement and confidence. It shows where the internal state adds information and where it does not. As a corollary, we show that the probe directions are readouts rather than mechanisms, and that the mean-difference direction is a steerable abstention gate.
  \item \textbf{Generalisation} to Llama and Mistral, to HotpotQA \cite{yang2018hotpotqa}, to RAGBench \cite{friel2024ragbench} across five domains, to a synthetic benchmark with paraphrased, partial and contradictory evidence, and to retrieved rather than curated evidence. We release the Evidence State Transition Benchmark (ESTB), which collects these sources under one labelling scheme, together with our code.
\end{itemize}
